% test-019.tex -- comando \spinterval
%
% Compilar:  lualatex-dev test-019.tex
% Validar:   show-pdf-tags --xml test-019.pdf | rnv-wrapp
%            verapdf --flavour ua2 --format text test-019.pdf
%
% Cada caso es un parrafo numerado, para poder referirse a el al escuchar
% el PDF. Los comentarios "doc:" de ESTE archivo son lecturas ESPERADAS por
% diseno, todavia no grabadas con NVDA: al verificar, reemplazarlas por las
% reales. Donde dice "..." el resto lo lee MathCAT. Los casos con \sqrt
% dependen de la regla de MathCAT para la raiz cuadrada.
%
% \spinterval recibe SIEMPRE el intervalo entre llaves: \spinterval{[a,b]}.
% Sin llaves, \spinterval[a,b] se leeria como las llaves de configuracion.
% Delimitadores: ( [ ] a la izquierda y ) ] [ a la derecha, en cualquier
% combinacion; ( ) ] [ excluyen el extremo y [ ] lo incluyen. El separador
% es una coma o, si hay decimales con coma, un punto y coma.
%
% Entradas que no van aqui porque detienen la compilacion: sin delimitadores,
% sin separador, un extremo vacio, mas de una coma sin punto y coma, llaves
% sin cerrar, un infinito incluido ([-\infty,3] o [2,\infty]) y el infinito
% del lado equivocado ([\infty,3] o [1,-\infty[).
\DocumentMetadata
  {
    lang=es-CL,
    pdfversion=2.0,
    pdfstandard=ua-2,
    tagging=on,
    tagging-setup={math/setup={mathml-SE}},
  }
\documentclass{article}
\usepackage[spanish]{babel}
\usepackage{unicode-math}
\usepackage{spintent}
\usepackage{hyperref}
\hypersetup
  {
    pdfdisplaydoctitle = true,
    pdftitle           = {test-019: spinterval},
  }
\pagestyle{empty}
\begin{document}

\section*{Los nueve pares de delimitadores}

% doc: «intervalo abierto de 1 a 5»
1. Paréntesis a los dos lados: $\spinterval{(1,5)}$.

% doc: «intervalo abierto de 1 a 5»
2. Paréntesis y corchete invertido: $\spinterval{(1,5[}$.

% doc: «intervalo abierto de 1 a 5»
3. Corchete invertido y paréntesis: $\spinterval{]1,5)}$.

% doc: «intervalo abierto de 1 a 5»
4. Corchetes invertidos, la notación chilena: $\spinterval{]1,5[}$.

% doc: «intervalo cerrado de 1 a 5»
5. Corchetes a los dos lados: $\spinterval{[1,5]}$.

% doc: «intervalo abierto en 1 y cerrado en 5»
6. Abierto a la izquierda con paréntesis y cerrado a la derecha: $\spinterval{(1,5]}$.

% doc: «intervalo abierto en 1 y cerrado en 5»
7. Abierto a la izquierda con corchete invertido y cerrado a la derecha: $\spinterval{]1,5]}$.

% doc: «intervalo cerrado en 1 y abierto en 5»
8. Cerrado a la izquierda y abierto a la derecha con paréntesis: $\spinterval{[1,5)}$.

% doc: «intervalo cerrado en 1 y abierto en 5»
9. Cerrado a la izquierda y abierto a la derecha con corchete invertido: $\spinterval{[1,5[}$.

\section*{Letras y subíndices}

% doc: «intervalo abierto de a a b»
10. Dos letras, abierto: $\spinterval{(a,b)}$.

% doc: «intervalo abierto de a a b»
11. Dos letras, abierto con corchetes invertidos: $\spinterval{]a,b[}$.

% doc: «intervalo cerrado de a a b»
12. Dos letras, cerrado: $\spinterval{[a,b]}$.

% doc: «intervalo abierto en a y cerrado en b»
13. Dos letras, abierto-cerrado: $\spinterval{]a,b]}$.

% doc: «intervalo cerrado en a y abierto en b»
14. Dos letras, cerrado-abierto: $\spinterval{[a,b[}$.

% doc: «intervalo cerrado de x sub uno a x sub dos»
15. Subíndices, cerrado: $\spinterval{[x_{1}, x_{2}]}$.

% doc: «intervalo abierto de x sub uno a x sub dos»
16. Subíndices, abierto: $\spinterval{]x_{1}, x_{2}[}$.

\section*{Signos y decimales}

% doc: «intervalo cerrado de menos 3 a 5»
17. Un extremo negativo: $\spinterval{[-3,5]}$.

% doc: «intervalo abierto de menos 2 a más 4»
18. Un signo en cada extremo: $\spinterval{]-2,+4[}$.

% doc: «intervalo cerrado de menos 4 a menos 1»
19. Dos extremos negativos: $\spinterval{[-4,-1]}$.

% doc: «intervalo cerrado de 3 coma 5 a 4 coma 2»
20. Decimales con coma y punto y coma como separador: $\spinterval{[3,5 ; 4,2]}$.

% doc: «intervalo abierto de 0 coma 5 a 1 coma 5»
21. Decimales entre 0 y 2, abierto: $\spinterval{]0,5 ; 1,5[}$.

% doc: «intervalo abierto en menos 1 coma 5 y cerrado en 2 coma 5»
22. Decimal negativo, abierto-cerrado: $\spinterval{]-1,5 ; 2,5]}$.

\section*{Infinito}

% doc: «intervalo abierto de menos infinito a 3»
23. Menos infinito a la izquierda, abierto: $\spinterval{]-\infty,3[}$.

% doc: «intervalo abierto en menos infinito y cerrado en 3»
24. Menos infinito a la izquierda, cerrado en la derecha: $\spinterval{]-\infty,3]}$.

% doc: «intervalo abierto de 2 a infinito»
25. Infinito a la derecha, abierto: $\spinterval{]2,\infty[}$.

% doc: «intervalo cerrado en 2 y abierto en infinito»
26. Infinito a la derecha, cerrado en la izquierda: $\spinterval{[2,\infty[}$.

% doc: «intervalo cerrado en 2 y abierto en más infinito»
27. Más infinito con signo explícito: $\spinterval{[2,+\infty[}$.

% doc: «intervalo abierto de menos infinito a infinito»
28. Los dos extremos infinitos: $\spinterval{]-\infty,\infty[}$.

% doc: «intervalo abierto de menos infinito a más infinito»
29. Los dos extremos infinitos con paréntesis y signo: $\spinterval{(-\infty,+\infty)}$.

\section*{Fracciones y tamaño de los delimitadores}

% doc: «intervalo cerrado de 1 mitad a 3»
30. Una fracción al comienzo, en línea: $\spinterval{[\frac{1}{2}, 3]}$.

% doc: «intervalo abierto de menos 1 mitad a 1 mitad»
31. Fracciones con signo, en línea: $\spinterval{]-\frac{1}{2}, \frac{1}{2}[}$.

% doc: «intervalo cerrado de 1 más 1 mitad a 3»
32. Una fracción que no está al comienzo: $\spinterval{[1+\frac{1}{2}, 3]}$.

% doc: «intervalo cerrado de 1 a 3 medios»
33. Una fracción con dfrac, en línea: $\spinterval{[1, \dfrac{3}{2}]}$.

% doc: «intervalo cerrado de raíz cuadrada de 1 mitad a 1»
34. Una fracción dentro de una raíz: $\spinterval{[\sqrt{\frac{1}{2}}, 1]}$.

% doc: «intervalo cerrado de 1 mitad a 3 medios»
35. Fracción y dfrac juntas: $\spinterval{[\frac{1}{2}, \dfrac{3}{2}]}$.

\section*{Raíces, potencias y subíndices}

% doc: «intervalo cerrado de raíz cuadrada de 2 a 3»
36. Una raíz en un extremo: $\spinterval{[\sqrt{2}, 3]}$.

% doc: «intervalo cerrado de menos raíz cuadrada de 2 a raíz cuadrada de 2»
37. Raíces con signo: $\spinterval{[-\sqrt{2}, \sqrt{2}]}$.

% doc: «intervalo cerrado de a cuadrado a b cuadrado»
38. Potencias: $\spinterval{[a^2, b^2]}$.

% doc: «intervalo abierto de x sub uno a menos raíz cuadrada de 2»
39. Subíndice y raíz con signo: $\spinterval{]x_{1}, -\sqrt{2}[}$.

\section*{La llave size}

% doc: «intervalo cerrado de 1 mitad a 3»
40. Con size auto y una fracción, en línea: $\spinterval[size=auto]{[\frac{1}{2}, 3]}$.

% doc: «intervalo cerrado de 1 mitad a 3»
41. Con size none y una fracción: $\spinterval[size=none]{[\frac{1}{2}, 3]}$.

% doc: «intervalo cerrado de 1 a 5»
42. Con size big: $\spinterval[size=big]{[1,5]}$.

% doc: «intervalo cerrado de 1 a 5»
43. Con size Big: $\spinterval[size=Big]{[1,5]}$.

% doc: «intervalo cerrado de 1 a 5»
44. Con size bigg: $\spinterval[size=bigg]{[1,5]}$.

% doc: «intervalo cerrado de 1 a 5»
45. Con size Bigg: $\spinterval[size=Bigg]{[1,5]}$.

% doc: «intervalo abierto de 1 a 5»
46. Con size Big y corchetes invertidos: $\spinterval[size=Big]{]1,5[}$.

% doc: «intervalo cerrado de raíz cuadrada de 2 a raíz cuadrada de 3»
47. Con size big y una raíz que lo necesita: $\spinterval[size=big]{[\sqrt{2}, \sqrt{3}]}$.

\section*{Las llaves prefix, concept, read-arg y silent}

% doc: «el intervalo abierto de 1 a 5»
48. Con prefix el: $\spinterval[prefix={el}]{(1,5)}$.

% doc: «un intervalo abierto en 1 y cerrado en 5»
49. Con prefix un y un mixto: $\spinterval[prefix={un}]{]1,5]}$.

% doc: «abierto de 1 a 5»
50. Con concept false, abierto: $\spinterval[concept=false]{(1,5)}$.

% doc: «abierto en 1 y cerrado en 5»
51. Con concept false, mixto: $\spinterval[concept=false]{]1,5]}$.

% doc: «abierto de 1 a 5»
52. Con concept false y prefix, que se ignora: $\spinterval[concept=false,prefix={el}]{(1,5)}$.

% doc: «rango»
53. Con read-arg de una palabra: $\spinterval[read-arg={rango}]{[1,5]}$.

% doc: «intervalo semiabierto»
54. Con read-arg de varias palabras: $\spinterval[read-arg={intervalo semiabierto}]{]1,5]}$.

% doc: «»
55. Con silent true: $\spinterval[silent=true]{[1,5]}$.

% doc: «»
56. Con silent true e infinito: $\spinterval[silent=true]{]-\infty,3[}$.

% doc: «»
57. Con silent true y una fracción: $\spinterval[silent=true]{[\frac{1}{2}, 3]}$.

% doc: «»
58. Con silent true y read-arg, gana silent: $\spinterval[silent=true,read-arg={rango}]{[1,5]}$.

\section*{El modo read-cmd mathcat}

% doc: «intervalo cerrado de 1 a 5»
59. Con mathcat, enteros: $\spinterval[read-cmd=mathcat]{[1,5]}$.

% doc: «intervalo cerrado de menos 3 a 5»
60. Con mathcat, un extremo negativo: $\spinterval[read-cmd=mathcat]{[-3,5]}$.

% doc: «intervalo cerrado de 3 coma 5 a 4 coma 2»
61. Con mathcat, decimales: $\spinterval[read-cmd=mathcat]{[3,5 ; 4,2]}$.

% doc: «intervalo abierto de menos 2 a más 4»
62. Con mathcat, signos y corchetes invertidos: $\spinterval[read-cmd=mathcat]{]-2,+4[}$.

% doc: «intervalo abierto en menos infinito y cerrado en 3»
63. Con mathcat, infinito: $\spinterval[read-cmd=mathcat]{]-\infty,3]}$.

% doc: «intervalo cerrado de 1 mitad a 3»
64. Con mathcat y una fracción: $\spinterval[read-cmd=mathcat]{[\frac{1}{2}, 3]}$.

\section*{Dentro de fórmulas}

% doc: «x … intervalo abierto de 1 a 5»
65. Un elemento en un intervalo: $x \in \spinterval{]1,5[}$.

% doc: «I es igual a intervalo cerrado de a a b»
66. Un intervalo igualado a otro símbolo: $I = \spinterval{[a,b]}$.

% doc: «A es igual a intervalo abierto en menos infinito y cerrado en 0»
67. Un conjunto igual a un intervalo con infinito: $A = \spinterval{]-\infty, 0]}$.

% doc: «intervalo abierto en menos infinito y cerrado en 2 unión intervalo cerrado en 3 y abierto en más infinito»
68. Unión de dos intervalos: $\spinterval{]-\infty, 2]} \cup \spinterval{[3, +\infty[}$.

% doc: «intervalo cerrado de 1 a 4 intersección intervalo abierto en 2 y cerrado en 6»
69. Intersección de dos intervalos: $\spinterval{[1,4]} \cap \spinterval{]2,6]}$.

% doc: «intervalo abierto de 1 a 5 … intervalo cerrado de 0 a 6»
70. Un intervalo contenido en otro: $\spinterval{]1,5[} \subset \spinterval{[0,6]}$.

% doc: «x … intervalo cerrado de 1 mitad a 3»
71. Un elemento en un intervalo con fracción: $x \in \spinterval{[\frac{1}{2}, 3]}$.

\section*{En fórmula destacada}

% doc: «x … intervalo cerrado de 1 mitad a 3»
72. Un elemento en un intervalo con fracción:
\[ x \in \spinterval{[\frac{1}{2}, 3]} \].

% doc: «intervalo abierto en menos infinito y cerrado en 2 unión intervalo cerrado en 3 y abierto en más infinito»
73. Unión de intervalos:
\[ \spinterval{]-\infty, 2]} \cup \spinterval{[3, +\infty[} \].

% doc: «intervalo cerrado de 1 a 3 medios»
74. Un intervalo con dfrac:
\[ \spinterval{[1, \dfrac{3}{2}]} \].

% doc: «A es igual a intervalo abierto de 1 mitad a 3 medios»
75. Un conjunto igual a un intervalo con fracciones:
\[ A = \spinterval{]\frac{1}{2}, \frac{3}{2}[} \].

\section*{Con setspintent}

\setspintent{\spinterval}{prefix={el}}%

% doc: «el intervalo abierto de 1 a 5»
76. Con setspintent y prefix el: $\spinterval{(1,5)}$.

% doc: «el intervalo cerrado de 2 a 6»
77. Otro intervalo con el mismo ajuste: $\spinterval{[2,6]}$.

\setspintent{\spinterval}{prefix={}}%

% doc: «intervalo abierto de 1 a 5»
78. Tras quitar el prefix: $\spinterval{(1,5)}$.

\setspintent{\spinterval}{size=Big}%

% doc: «intervalo cerrado de 1 a 5»
79. Con setspintent y size Big: $\spinterval{[1,5]}$.

\setspintent{\spinterval}{size=auto}%

% doc: «intervalo cerrado de 1 a 5»
80. Tras volver a size auto: $\spinterval{[1,5]}$.

\end{document}
